\section{Введение}

Многомерные данные перестали быть принадлежностью специализированных систем.
Геометрические и географические объекты, диапазоны значений, полнотекстовые
сигнатуры и векторные представления обрабатываются универсальными СУБД наравне
со скалярными типами, и универсальным механизмом доступа к ним стало обобщённое
дерево поиска~\cite{hellerstein1995generalized}: ядро реализует страничную
организацию, спуск, расщепление, конкурентный доступ и журналирование, а
предметная часть выносится в класс операторов.

Теория качества таких деревьев развита: со времён R-дерева~\cite{guttman1984r}
известно, что стоимость поиска определяется перекрытием ключей узлов и покрытой
ими площадью, и предложены критерии разбиения узла, эти величины
улучшающие~\cite{beckmann1990rstar,beckmann2009rrstar}. Аналитическая модель
предсказания стоимости запроса к R-дереву построена Теодоридисом и Селлисом
\cite{theodoridis1996model}: число просмотренных узлов выражается через
плотность данных и размеры ограничивающих прямоугольников, оцениваемые по
выборке. Свод результатов этого направления дан
в~\cite{manolopoulos2006rtrees}, обзор многомерных методов доступа~---
в~\cite{gaede1998multidimensional}.

Модель~\cite{theodoridis1996model} была основанием и для более ранних работ
авторов, где оценивалась эффективность агрегирующих запросов к аналитическим
данным, отображённым в пространственные~\cite{borodin2009rtree,
borodin2011estimation, borodin2012monograph}; сопоставление с реальными
измерениями на средствах бизнес-аналитики приведено в~\cite{borodin2010bi}.
Настоящая работа продолжает эту линию, но меняет предмет: там оценивалась
стоимость запроса к дереву заданной формы, здесь~--- зависимость этой формы от
решений, принимаемых реализацией.

Общее свойство перечисленных моделей~--- они выражены в терминах свойств
разбиения пространства (плотность, площадь, перекрытие) и не связаны с
параметрами страничной организации узла, тогда как в реализации управлять можно
именно последними. Разработчик СУБД не может
изменить перекрытие непосредственно: он может изменить представление ключа,
способ выбора точек разбиения, порядок обхода страниц при построении и
сопровождении. Какое из этих действий и на сколько изменит наблюдаемую
стоимость запроса~--- вопрос, на который существующие модели не отвечают.

Практическое следствие этого разрыва состоит в том, что методы построения
индекса сравниваются по времени построения. Такое сравнение систематически
вводит в заблуждение: способ построения влияет на форму дерева, то есть на
стоимость всех последующих запросов, и выигрыш во времени построения может
быть с избытком оплачен.

Цель работы~--- построить модели стоимости построения, поиска и сопровождения
обобщённого дерева поиска, выраженные через параметры, которыми располагает
реализация, и проверить их измерением на промышленной СУБД.

\section{Постановка и обозначения}

Рассматривается сбалансированное дерево, узел которого занимает страницу.
Элемент внутреннего узла~--- пара «ключ, нисходящая ссылка»; ключ покрывает
ключи всех элементов поддерева. Поиск спускается во все поддеревья, ключи
которых могут содержать ответ, поэтому путей спуска может быть несколько, и
стоимость поиска определяется их числом, а не высотой дерева.

Класс операторов предоставляет операции \code{consistent} (может ли поддерево
содержать ответ), \code{union} (ключ, покрывающий набор ключей), \code{penalty}
(штраф за расширение ключа) и \code{picksplit} (разбиение переполненного узла).
Порядка на множестве ключей интерфейс не содержит.

\begin{table}[h]
\centering
\caption{Обозначения}
\begin{tabular}{ll}
\toprule
$N$ & число индексируемых записей \\
$B$ & размер страницы \\
$h$ & размер служебного заголовка страницы \\
$k$ & средний размер ключа \\
$p$ & размер ссылки \\
$\sigma$ & служебные данные, сопровождающие элемент в узле \\
$f$ & ветвистость узла \\
$H$ & высота дерева \\
$\varphi$ & коэффициент заполнения страницы \\
$D$ & размерность данных \\
$M$ & память, доступная сортировке \\
$c_r$, $c_s$ & стоимость произвольного и последовательного чтения страницы \\
$c_c$ & стоимость одного вычисления \code{consistent} \\
\bottomrule
\end{tabular}
\end{table}

\section{Ветвистость узла}

Ветвистость определяется размером страницы и размером хранимого элемента:

\begin{equation}
\label{eq:fanout}
f = \left\lfloor \frac{B-h}{k+p+\sigma} \right\rfloor ,
\qquad
H = \left\lceil \log_f \frac{N}{f} \right\rceil + 1,
\qquad
L = \left\lceil \frac{N}{f\varphi} \right\rceil .
\end{equation}

Слагаемое $\sigma$ в теоретических работах обычно опускают. В реализации оно не
равно нулю: ключ может храниться в преобразованном виде, вместе с ним могут
храниться дополнительные атрибуты, а часть интерфейса класса операторов
вынуждает хранить избыточные представления. Именно $\sigma$ и $k$~--- те
величины, на которые реализация воздействует непосредственно.

Из~\eqref{eq:fanout} видно, что ветвистость влияет на высоту логарифмически, а
на число страниц~--- обратно пропорционально. Поэтому эффект от её роста
проявляется не столько в сокращении длины пути спуска, сколько в сокращении
объёма индекса и стоимости операций, обходящих индекс целиком: построения,
сборки мусора, верификации.

\section{Стоимость поиска и перекрытие ключей}

Определим перекрытие уровня $l$ как математическое ожидание числа узлов этого
уровня, ключи которых удовлетворяют условию поиска:

\begin{equation}
\label{eq:omega}
\omega_l = \mathbb{E}\bigl|\{\,u \in U_l : \code{consistent}(key(u), Q)\,\}\bigr| .
\end{equation}

Для дерева без перекрытия $\omega_l = 1$ на всех уровнях. Стоимость точечного
запроса составляет

\begin{equation}
\label{eq:point}
P_{\text{point}} = \sum_{l=0}^{H-1} \omega_l \, c_r .
\end{equation}

Узел посещается, только если посещён его родитель, что подсказывает вывод о
накоплении перекрытия вниз по дереву. \emph{Этот вывод неверен.} Измерение
(§\ref{sec:exp}) показывает, что $\omega_l$ выходит на характерное для способа
построения значение $\omega_\infty$ уже на первом-втором уровне под корнем и
далее с глубиной не растёт; верхние уровни составляют исключение лишь потому,
что там $\omega_l$ ограничено числом узлов уровня. Следовательно

\begin{equation}
\label{eq:steady}
P_{\text{point}} \approx \omega_\infty \, H \, c_r ,
\end{equation}

и качество разбиения входит в стоимость поиска постоянным множителем на каждом
шаге спуска, а не накапливающимся по уровням.

Различие между~\eqref{eq:point} и~\eqref{eq:steady} существенно для практики.
Мультипликативное накопление означало бы, что ухудшение разбиения тем опаснее,
чем больше данных; постоянный множитель означает, что относительный проигрыш от
плохого разбиения не зависит от объёма. Второе и наблюдается.

\subsection{Процессорная составляющая}

Формула~\eqref{eq:point} учитывает только обращения к страницам. При данных,
помещающихся в оперативной памяти, доминирует другая величина: чтобы отобрать
поддеревья для спуска, ядро вычисляет \code{consistent} для каждого из $f$
элементов узла. С учётом этого

\begin{equation}
\label{eq:cpu}
T = \sum_{l=0}^{H-1} \omega_l \left( c_r + f\,c_c \right) .
\end{equation}

Из~\eqref{eq:cpu} следует ограничение, не видное в~\eqref{eq:point}: рост
ветвистости сокращает число обращений к страницам, но линейно увеличивает
процессорную стоимость просмотра узла. Осмысленное увеличение размера страницы
поэтому ограничено, и снятие этого ограничения требует внутристраничной
структуры, позволяющей отбирать кандидатов не за $O(f)$.

\section{Влияние размерности}

Пусть данные равномерно распределены в единичном $D$-мерном кубе, а ключ
узла~--- ограничивающий параллелепипед со стороной $a_l$. Вероятность того, что
случайный точечный запрос попадёт в ключ, равна $a_l^D$, откуда

\begin{equation}
\label{eq:curse}
\omega_l \approx |U_l| \, a_l^{D}, \qquad
a_l \approx \left(\frac{f^{\,l+1}}{N}\right)^{1/D} .
\end{equation}

Показатель $1/D$ означает, что при больших $D$ сторона ключа близка к единице
уже на средних уровнях, и малое отклонение разбиения от идеального даёт кратный
рост $\omega_l$. Отсюда объяснение практического наблюдения: приёмы, безопасные
для одномерного B-дерева, для многомерного дерева опасны, поскольку там
качество разбиения входит в стоимость с показателем, растущим по размерности.

\section{Стоимость построения}

Построение последовательной вставкой спускается к листу для каждой записи:

\begin{equation}
\label{eq:build-ins}
P_{\text{ins}} = N\left(H + \gamma\right) c_r ,
\end{equation}

где $\gamma$~--- среднее число дополнительных обращений на расщепление и
обновление ключей родителей. Обращения произвольные: порядок вставки не связан
с физическим расположением страниц.

Если для класса операторов задан порядок, сохраняющий локальность, набор данных
можно отсортировать и заполнять страницы подряд снизу вверх:

\begin{equation}
\label{eq:build-sort}
P_{\text{sort}} = \frac{2N}{B}\log_{M/B}\frac{N}{B}\, c_s
                + \frac{N}{f\varphi}\, c_s .
\end{equation}

Источников выигрыша три: замена произвольного доступа последовательным, отказ
от расщеплений и рост $\varphi$ с характерных для расщепления $0{,}7$ до
заданного коэффициента заполнения.

\subsection{Критерий сопоставления методов построения}

Сравнение~\eqref{eq:build-ins} и~\eqref{eq:build-sort} неполно, поскольку метод
построения влияет на $\omega_\infty$, то есть на стоимость всех последующих
запросов. Правильным основанием для сравнения является

\begin{equation}
\label{eq:criterion}
P_{\text{build}} + n\,P_{\text{point}}
  \;=\; P_{\text{build}} + n\,\omega_\infty H c_r ,
\end{equation}

где $n$~--- ожидаемое число запросов за время жизни индекса. Из~\eqref{eq:criterion}
следует, что метод построения, ускоряющий сборку в $\alpha$ раз ценой роста
$\omega_\infty$ в $\beta$ раз, оправдан лишь при

\begin{equation}
\label{eq:break-even}
n < \frac{P_{\text{build}}\,(1 - 1/\alpha)}{(\beta - 1)\,\omega_\infty H c_r} ,
\end{equation}

то есть при достаточно коротком времени жизни индекса или достаточно редких
запросах. Для индекса, строящегося однократно и обслуживающего миллионы
запросов, правая часть~\eqref{eq:break-even} пренебрежимо мала, и любой рост
$\omega_\infty$ неприемлем независимо от выигрыша в сборке.

\section{Стоимость сопровождения}

Сборка мусора обходит все страницы индекса. При обходе в логическом порядке
последовательность номеров страниц произвольна, при обходе в физическом
порядке~--- возрастающая:

\begin{equation}
\label{eq:vac}
P^{\text{log}}_{\text{vac}} = (L+I)\,c_r ,
\qquad
P^{\text{phys}}_{\text{vac}} = (L+I)\,c_s + \beta\,c_r ,
\end{equation}

где $I$~--- число внутренних страниц, $\beta$~--- число страниц, потребовавших
повторной обработки из-за конкурентных вставок. Выигрыш пропорционален
$(L+I)(c_r-c_s)$ и, следовательно, объёму индекса~--- то есть обратно
пропорционален ветвистости, что связывает~\eqref{eq:vac} с~\eqref{eq:fanout}.

\section{Экспериментальная проверка}
\label{sec:exp}

Измерения выполнены на СУБД PostgreSQL, 16-ядерная машина, буферный кэш 8~ГиБ.
Первичная величина~--- число обращений к страницам индекса на запрос, полученное
из плана выполнения. Она равна $\sum_l \omega_l$, не зависит от носителя и
воспроизводится при фиксированном зерне генератора точно, в отличие от
измерений времени.

\subsection{Перекрытие не накапливается с глубиной}

Таблица~\ref{tab:levels} приводит $\omega_l$ по уровням для десяти миллионов
равномерно распределённых точек при трёх способах построения.

\begin{table}[h]
\centering
\caption{Перекрытие по уровням, $N=10^7$, 2000 зондов}
\label{tab:levels}
\begin{tabular}{lrrrrrr}
\toprule
& \multicolumn{2}{c}{вставка} & \multicolumn{2}{c}{по заполнению} & \multicolumn{2}{c}{\code{picksplit}} \\
\cmidrule(lr){2-3}\cmidrule(lr){4-5}\cmidrule(lr){6-7}
Уровень & узлов & $\omega_l$ & узлов & $\omega_l$ & узлов & $\omega_l$ \\
\midrule
0 & 1 & 1,00 & 1 & 1,00 & 1 & 1,00 \\
1 & 8 & 1,05 & 3 & 1,78 & 5 & 1,99 \\
2 & 823 & 1,13 & 363 & 4,24 & 601 & 2,47 \\
3 & 90\,076 & 1,06 & 60\,241 & 3,95 & 80\,737 & 1,35 \\
\bottomrule
\end{tabular}
\end{table}

Каждый способ выходит на своё значение в пределах одного-двух уровней от корня
и далее его удерживает: $\omega_\infty \approx 1{,}05$ для построения вставкой,
$\approx 4{,}0$ для разбиения по заполнению страницы, $\approx 1{,}3$ для
разбиения функцией класса операторов. Это и есть эмпирическое содержание
формулы~\eqref{eq:steady}.

Следствие проверяется независимо: при увеличении объёма данных на порядок
высота дерева растёт, а отношение стоимостей между способами построения должно
сохраниться. Измерено 12,67 против 6,86 обращений при $10^7$ и 15,20 против
7,91 при $1{,}38\cdot10^8$ точках реальных данных, то есть отношение 1,85 и
1,92 соответственно.

\subsection{Перекрытие, а не ветвистость}

Стоимость запроса в измерениях монотонно убывает с ростом числа страниц
индекса, что допускает истолкование через ветвистость: менее ветвистое дерево
имеет меньше элементов в узле, и классическая оценка даёт оптимальную
ветвистость около $e$. Для проверки коэффициент заполнения при разбиении по
заполнению страницы понижался до сравнивания числа страниц с тем, что даёт
разбиение через \code{picksplit}.

\begin{table}[h]
\centering
\caption{Проверка истолкования через ветвистость, $N=10^7$}
\label{tab:control}
\begin{tabular}{llrr}
\toprule
разбиение & $\varphi$ & страниц & обращений \\
\midrule
\code{picksplit} & по умолчанию & 81\,326 & \textbf{6,86} \\
по заполнению & по умолчанию & 60\,608 & 12,67 \\
по заполнению & 65 & 84\,036 & 13,64 \\
по заполнению & 50 & 109\,892 & 14,49 \\
\bottomrule
\end{tabular}
\end{table}

При практически равном числе страниц (84\,036 против 81\,326) разбиение по
заполнению требует вдвое больше обращений. Снижение коэффициента заполнения не
улучшает стоимость запроса, а слегка ухудшает её. Истолкование через
ветвистость не подтверждается: в стоимость входит качество разбиения, то есть
$\omega_\infty$, а не число страниц.

\subsection{Процессорная составляющая}

Профилирование серверного процесса на нагрузке из пространственных соединений,
когда индекс и таблица целиком находятся в буферном кэше, даёт следующее
распределение: перебор элементов страницы 22,4\,\%, вычисление предиката
14,1\,\%, поиск страницы в таблице буферов 12,8\,\%, закрепление страницы
5,7\,\%, вызов предиката через диспетчер функций 4,7\,\%. Работа с элементами
страницы составляет в сумме 46,8\,\%.

Это подтверждает~\eqref{eq:cpu} и уточняет её: величина $c_c$ складывается не
столько из предметных вычислений, сколько из накладных расходов обхода. Прямая
проверка следствия: реализация внутристраничной структуры, позволяющей
пропускать группы элементов, ускорила соединение по индексу в 1,41 раза при
неизменном числе обращений к страницам, то есть выигрыш получен целиком за счёт
слагаемого $f c_c$.

\section{Границы применимости моделей}

Модели проверены не полностью, и это следует назвать явно.

\emph{Зависимость от размерности}~\eqref{eq:curse} получена аналитически и
измерением не подтверждена. Проверка требует класса операторов над данными
переменной размерности; для него недоступно сортированное построение, поэтому
измерение возможно только для дерева, построенного вставкой.

\emph{Стоимость сопровождения}~\eqref{eq:vac} измерением не проверялась.
Разделение произвольного и последовательного чтения требует носителя с
известными характеристиками, тогда как измерения выполнялись на виртуальной
машине, где отношение $c_r/c_s$ определяется гипервизором, а не устройством.

\emph{Установившееся перекрытие} измерено для одного класса операторов при
равномерном и кластеризованном распределениях и для одного набора реальных
данных. Утверждение о том, что $\omega_l$ выходит на постоянное значение,
опирается на деревья высотой три и четыре; для более высоких деревьев оно не
проверялось.

\emph{Величина $c_c$} в~\eqref{eq:cpu} измерена профилированием для одного
класса операторов. Для классов с существенно более дорогим предикатом её доля
должна быть выше, что согласуется с моделью, но не измерено.

\section{Заключение}

Построены модели стоимости построения, поиска и сопровождения обобщённого
дерева поиска, выраженные через параметры страничной организации узла.
Установлено, что перекрытие ключей не накапливается с глубиной дерева, а
выходит на характерное для способа построения установившееся значение, и потому
входит в стоимость поиска постоянным множителем на каждом шаге спуска.

Из моделей выделены три величины, поддающиеся воздействию без изменения класса
операторов: ветвистость узла, установившееся перекрытие и порядок обращения к
страницам. Показано, что при данных, помещающихся в оперативной памяти,
ветвистость входит в стоимость дважды и с противоположными знаками:
сокращая число обращений к страницам и увеличивая процессорную стоимость
просмотра узла.

Предложен критерий сопоставления методов построения индекса, учитывающий
влияние способа построения на стоимость последующих запросов, и получено
условие~\eqref{eq:break-even}, при котором ускорение построения ценой ухудшения
дерева оправдано.

Все утверждения проверены измерением на промышленной СУБД; первичной величиной
выбрано число обращений к страницам индекса как не зависящее от стенда и
воспроизводимое точно.
